\clearpage\begingroup\NoSourceHead
\makeatletter\tagsleft@true\let\veqno\@@leqno\makeatother
\renewcommand{\thefootnote}{\arabic{footnote})}\setcounter{footnote}{28}
\ArithmeticHead{18}{L}{BOLZA, GAUSS UND DIE VARIATIONSRECHNUNG.}
\noindent\hspace*{\parindent}9. \textsc{Lagrange} hatte die Formel (16) für $\delta p$, $\delta q$ und damit auch das Schlussresultat
(22) nur unter der dem Problem selbst durchaus fernliegenden beschränkenden
Voraussetzung (15) abgeleitet, und überdies hat er dabei die
fragwürdige Operation der »formalen Variation« angewandt.

Die hier noch vorliegende Lücke wurde bald darauf (1816) von \textsc{Poisson}\footnote{\textit{Sur le calcul des variations, relativement aux intégrales multiples}, Bulletin des Sciences, par la Société Philomatique de Paris, Année 1816, S. 82—86; wiedergegeben in \textsc{Lacroix}, \textit{Traité etc.}, Bd. III (1819), S. 717—721.}
ausgefüllt. Er leitet zunächst mit \textsc{Lagrange} die Formel (19) mittels »formaler
Variation« ab. Dann aber führt er an Stelle von $x,y$ zwei Hilfsvariable $u,v$
als unabhängige Variable ein. Dadurch gehen dann auch $z$ und die partiellen
Ableitungen
\[
p=\frac{\partial z}{\partial x},\quad q=\frac{\partial z}{\partial y}
\]
in Funktionen von $u,v$ über. Jetzt werden die Funktionen $x,y,z$ von $u,v$
variiert, während $u,v$ unvariiert bleiben. Dann ergeben sich die Variationen
von $\delta p$, $\delta q$ zunächst als Funktionen von $u,v$ und weiterhin, indem man von
den Variabeln $u,v$ zu den alten Variabeln $x,y$ zurückkehrt, als Funktionen
von $x,y$. Am einfachsten geschieht dies, indem man direkt $x=u$, $y=v$
setzt, was keine Beschränkung der Allgemeinheit ist.

Als Resultat erhält \textsc{Poisson} die \textsc{Lagrange}schen Ausdrücke (16) für $\delta p$, $\delta q$,
\par\clearpage\SourceLeaf{476}
\ArithmeticHead{19}{R}{GESCHICHTE DER EXTREMA VON DOPPELINTEGRALEN VOR GAUSS.}
\noindent von denen nunmehr bewiesen ist, dass sie nicht nur unter der beschränkenden
Voraussetzung (15), sondern \textls[100]{ganz allgemein gültig} sind.

Zugleich ergab sich eine Bestätigung der von \textsc{Lagrange} durch eine geometrische
Infinitesimalbetrachtung erhaltenen Formel (20) für die Variation
des Flächenelementes, das jetzt die Form
\[
dx\,dy=\left(\frac{\partial x}{\partial u}\frac{\partial y}{\partial v}-\frac{\partial x}{\partial v}\frac{\partial y}{\partial u}\right)du\,dv
\]
annimmt.

Aber die grundsätzliche Bedeutung der Arbeit von \textsc{Poisson} reicht weit
über die angeführten Verbesserungen der Untersuchungen von \textsc{Lagrange} hinaus,
insofern sie die Mittel liefert, die Berechnung der Variation des Doppelintegrals
von der anfechtbaren Operation der »formalen Variation« ganz frei
zu machen. Zwar hat \textsc{Poisson} in dieser ersten Arbeit bei Ableitung der
Gleichung (19) noch keinen Gebrauch davon gemacht. Aber vorgreifend wollen
wir doch schon hier erwähnen, dass er in seiner spätern grossen Arbeit »\textit{Mémoire
sur le calcul des variations}«\footnote{Mémoires de l'Académie royale des Sciences de l'Institut de France, Bd. XII (1833), S. 287—293. Datiert ist die Arbeit vom 10. November 1831.} (1831) auch diesen letzten Rest der »formalen
Variation« beseitigt hat, indem er die unabhängigen Variabeln $u,v$
gleich in das Integral $J$ einführt. Dasselbe geht dann über in
\begin{equation*}
J=\iint V\left(\frac{\partial x}{\partial u}\frac{\partial y}{\partial v}-\frac{\partial x}{\partial v}\frac{\partial y}{\partial u}\right)du\,dv,\tag{24}
\end{equation*}
erstreckt über einen gewissen Bereich im Gebiet der Variabeln $u,v$. Nunmehr
werden $x,y,z$ durch $x+\delta x$, $y+\delta y$, $z+\delta z$ ersetzt, wobei $\delta x$, $\delta y$, $\delta z$ willkürliche
unendlich kleine Funktionen von $u,v$ sind, während $u,v$ auf denselben
Bereich beschränkt werden wie im Integral $J$. Auf diese Weise kann
man von einem gegebenen Flächenstück zu einem beliebigen benachbarten
übergehen, ohne den Integrationsbereich zu ändern. Bei der Variation von
$J$ hat man daher jetzt weder die Grenzen noch die unabhängigen Variabeln
zu variieren, sondern einzig die Funktionen $x,y,z$, was keinerlei begriffliche
Schwierigkeiten mehr bietet.

Damit waren endlich die mit der Variation der unabhängigen Variabeln
verbundenen Schwierigkeiten und Unklarheiten aus dem Wege geräumt, freilich
\par\smallskip\noindent\hfill{\small 3\textsuperscript{*}}\par
\clearpage\SourceLeaf{477}
\ArithmeticHead{20}{L}{BOLZA, GAUSS UND DIE VARIATIONSRECHNUNG.}
\noindent nur dadurch, dass überhaupt der ganze Begriff der Variation der unabhängigen
Variabeln beseitigt wurde.

Auf die mit der partiellen Integration verbundenen Schwierigkeiten ist
auch \textsc{Poisson} in seiner Arbeit von 1816 nicht eingegangen. Er begnügt sich
mit der \textsc{Euler}schen Gleichung (5) und das Problem der Randbedingung blieb
ungelöst.

Ein schlagendes Licht auf die hier noch vorliegende Lücke wirft die
Diskussion, die sich an die von \textsc{Gergonne} in demselben Jahr 1816 gestellte
Aufgabe geknüpft hat\footnote{Annales de Mathématiques pures et appliquées, Bd. VII (1816/7), S. 99, 148—155, 283—287, 284/5 Fussnote.}: Einen Würfel in der Art in zwei Teile zu schneiden,
dass die Schnittfläche durch zwei nicht parallele Diagonalen zweier gegenüberliegender
Würfelflächen geht und einen möglichst kleinen Flächeninhalt besitzt.
Als Lösung gab \textsc{Tédénat}\begingroup\renewcommand{\thefootnote}{31a)}\footnote{Ebenda, S. 283.}\endgroup\setcounter{footnote}{31}\ eine Schraubenfläche an. \textsc{Gergonne} äussert
Bedenken gegen die Richtigkeit dieser Lösung. Aber zum direkten Nachweis
ihrer Unrichtigkeit fehlten ihm noch die Hilfsmittel, da ihm selbst für dieses
denkbar einfachste Beispiel, das sich ohne Variation der Grenzen erledigen
liess\footnote{Vgl. die Bemerkungen von \textsc{H. A. Schwarz} zur \textsc{Gergonne}schen Aufgabe, Monatsberichte der Berliner Akademie, 1872, S. 3 ff., Gesammelte Abhandlungen, Bd. I, S. 126—129.}, die Randbedingung unbekannt war, die hier darin besteht, dass die
gesuchte Fläche in den nicht vorgegebenen Teilen ihrer Begrenzung die Wände
des Würfels senkrecht schneiden muss.

\begin{center}
\textls[100]{II. Abschnitt: Die »Principia generalia theoriae figurae fluidorum}\par
\textls[100]{in statu aequilibrii«.}
\end{center}

10. Die beiden Fundamentalformeln (9) und (13) von \textsc{Ohm} einerseits und
die von \textsc{Poisson} bewiesene \textsc{Lagrange}sche Formel (22) andererseits bezeichnen
die äussersten Punkte, welche die Theorie der Variation mehrfacher Integrale
erreicht hatte, als \textsc{Gauss} mit seiner Arbeit \textit{Principia generalia theoriae figurae fluidorum
in statu aequilibrii}\footnote{Der Göttinger Gesellschaft der Wissenschaften vorgelegt am 28. September 1829, veröffentlicht als Sonderausgabe 1830 und in den Commentationes societatis regiae scientiarum Gottingensis recentiores, Vol. VII (1832), Class. math. S. 39—88; Werke, Bd. V, S. 29—77.} in die Variationsrechnung eingriff. In dieser Arbeit
leitet \textsc{Gauss} aus dem Prinzip der virtuellen Verrückungen für das Gleichgewicht
\par\clearpage\SourceLeaf{478}
\ArithmeticHead{21}{R}{GAUSS’ PRINCIPIA GENERALIA THEORIAE FIGURAE FLUIDORUM.}
\noindent einer homogenen, inkompressibeln Flüssigkeit, deren Bewegungsfreiheit
durch einen unbeweglichen starren Körper (Gefäss) beschränkt ist, und auf
deren Teilchen, ausser der Schwerkraft, einerseits die gegenseitige Anziehung
der Flüssigkeitsteilchen, andererseits die Anziehung der Teilchen des Gefässes
wirken, die folgende Bedingung ab:

Es bezeichne $s$ das Volumen der Flüssigkeit, $T$ den Inhalt desjenigen
Teiles der Oberfläche der Flüssigkeit, der das Gefäss berührt, $U$ den Inhalt
des andern (freien) Teiles der Oberfläche der Flüssigkeit, $z$ den vertikalen
Abstand eines Punktes der Flüssigkeit über einer festen horizontalen Ebene,
endlich $\alpha$, $\beta$ zwei positive physikalische Konstanten. Alsdann muss im Zustand
des Gleichgewichts der Ausdruck
\begin{equation*}
W=\int z\,ds+(\alpha^2-2\beta^2)T+\alpha^2U\tag{25}
\end{equation*}
ein Minimum sein in Beziehung auf die Gesamtheit aller mit der Gestalt des
Gefässes vereinbaren unendlich kleinen Formveränderungen der Flüssigkeit,
bei denen ihr Volumen unverändert bleibt.

Nimmt man mit \textsc{Gauss} an, dass die freie Oberfläche der Flüssigkeit, die
ebenso wie ihr Inhalt mit $U$ bezeichnet wird, sich in rechtwinkligen Koordinaten
in der Form $z=z(x,y)$ darstellen lässt, und macht man dieselbe Annahme
für die Gefässwand und schreibt ihre Gleichung: $z=g(x,y)$, so lässt
sich $W$ durch die folgende Summe von Doppelintegralen ausdrücken
\begin{equation*}
\begin{aligned}
W={}&\tfrac12\iint[z^2-g^2(x,y)]\,dx\,dy\\
&+(\alpha^2-2\beta^2)\iint\sqrt{1+g_x^2+g_y^2}\,dx\,dy\\
&+\alpha^2\iint\sqrt{1+p^2+q^2}\,dx\,dy,
\end{aligned}\tag{26}
\end{equation*}
während das vorgegebene Volumen ausgedrückt wird durch
\begin{equation*}
s=\iint[z-g(x,y)]\,dx\,dy.\tag{27}
\end{equation*}
Der Integrationsbereich ist dabei die Projektion $\mathfrak{A}$ der freien Oberfläche $U$
auf die $x,y$-Ebene.

Es handelt sich also um ein \textls[100]{zweidimensionales isoperimetrisches
Problem mit variablem Integrationsbereich}, und \textsc{Gauss} hätte, wenn
er von den Resultaten seiner Vorgänger hätte Gebrauch machen wollen, (vorausgesetzt
dass er dieselben gekannt hat), für die Berechnung von $\delta W$ und $\delta s$
\par\clearpage\SourceLeaf{479}
\ArithmeticHead{22}{L}{BOLZA, GAUSS UND DIE VARIATIONSRECHNUNG.}
\noindent von den fertig vorliegenden Formeln (9) und (13) von \textsc{Ohm} oder von der Formel
(22) von \textsc{Lagrange} ausgehen können.

\textsc{Gauss} hat diesen Weg jedoch nicht eingeschlagen. Die Untersuchungen
von \textsc{Ohm} scheinen ihm nicht bekannt gewesen zu sein, wie man wohl aus
dem folgenden Satz (Art. 20) schliessen kann: »Sed quum calculus variationum
integralium duplicium pro casu, ubi etiam limites tamquam variabiles spectari
debent, hactenus parum excultus sit, hanc disquisitionem subtilem paullo profundius
petere oportet«, sowie auch aus dem folgenden Passus, in dem \textsc{Gauss}
der Methode der Variation der unabhängigen Variabeln vor der Methode der
Variation der Grenzen den Vorzug gibt (Art. 21): »Si sufficeret, tales tantummodo
mutationes considerare, pro quibus limes $P$ semper invariatus, vel saltem
in eadem superficie verticali maneret, manifesto soli coordinatae tertiae $z$
variationem inducere oporteret, quo pacto problema longe facilius evaderet;
sed quum problema maxima generalitate nobis ventilandum sit, in tali investigationis
modo consideratio variabilitatis limitum in ambages incommodas concinnitatemque
turbantes perduceret; quamobrem praestabit, statim ab initio
omnes tres coordinatas variationi subiicere«.

Aber auch von den \textsc{Lagrange}schen Formeln (21) und (22) macht \textsc{Gauss}
keinen Gebrauch. Vielmehr hat er es vorgezogen, unabhängig von allen vorangegangenen
Untersuchungen über die Variation von Doppelintegralen und unter
Beiseitesetzung des ganzen Apparates des $\delta$-Algorithmus unmittelbar aus der
geometrischen Bedeutung der Integrale $U$, $T$, $s$, $\int z\,ds$ die Variation dieser
Grössen zu berechnen.

Sehen wir nun zu, wie \textsc{Gauss} im einzelnen das von ihm gestellte Variationsproblem
löst. Entsprechend den Hauptschwierigkeiten, denen wir in der Geschichte
der Extrema von Doppelintegralen vor \textsc{Gauss} begegnet sind, werden
wir dabei drei Etappen unterscheiden:
\par\noindent\hspace*{\parindent}A) Die Aufstellung der ersten Variation,
\par\noindent\hspace*{\parindent}B) Die Umformung der ersten Variation mittels partieller Integration,
\par\noindent\hspace*{\parindent}C) Die Ableitung der partiellen Differentialgleichung und der Randbedingung.
\par\noindent\hspace*{\parindent}Wir bemerken noch, dass \textsc{Gauss} auf eine Untersuchung der zweiten Variation
nicht eingegangen ist\footnote{Mit einer einzigen Ausnahme: In Art. 19 behandelt \textsc{Gauss}, ehe er das allgemeine Problem in Angriff nimmt, den speziellen Fall von kommunizierenden Röhren mit vertikalen zylindrischen Schenkeln.
Durch Betrachtung einer speziellen Art von Variationen, bei denen die freien Oberflächenstücke der Flüssigkeit zu sich selbst parallel in vertikaler Richtung verschoben werden, leitet er in sehr einfacher Weise das Gesetz für die Steighöhen ab. Hierbei geht er auch auf die Glieder zweiter Ordnung ein.}, dieselbe überhaupt nicht erwähnt.
\par\clearpage\endgroup
